\chapter{Boundary Conformal Field Theory}
\label{ch:boundary-conformal-field-theory}

\ConventionsEuclideanCFT

\section{Boundary CFT Data}

Boundary conformal field theory is not just a conformal field theory placed
near a wall.  The boundary condition is additional local QFT data: it
specifies boundary operators, bulk-to-boundary operator expansions, and
sewing maps for bordered surfaces.  In two dimensions this datum admits a
particularly sharp operator formulation because the conformal boundary
condition identifies the holomorphic and antiholomorphic stress tensors on
the boundary.

\begin{definition}[Oriented bosonic BCFT datum]
\label{def:oriented-bosonic-bcft-datum}
Fix a closed oriented two-dimensional CFT \(\mathcal C\).  An oriented
bosonic boundary conformal field theory based on \(\mathcal C\) consists of
the following data.
\begin{enumerate}
\item A set \(\mathsf B\) of boundary conditions.
\item For each ordered pair \(a,b\in\mathsf B\), a positive-energy Hilbert
space \(\Hilb_{ab}\) assigned to an interval whose left boundary condition is
\(a\) and whose right boundary condition is \(b\).
\item For each \(a\in\mathsf B\), a boundary operator algebra
\(\mathcal A_a\) acting on boundary segments with condition \(a\), and more
generally boundary-condition-changing operators in
\(\operatorname{Hom}(\Hilb_{ab},\Hilb_{ac})\).
\item Bulk-to-boundary expansion maps expressing a bulk local field
\(\mathcal O(z,\bar z)\), as \(z\) approaches a boundary segment with
condition \(a\), in the form
\begin{equation}
  \mathcal O(x+\ii y,x-\ii y)
  \sim
  \sum_{\beta}
  B_{\mathcal O}^{a,\beta}\,
  (2y)^{h_\beta-h-\bar h}
  \widehat{\mathcal O}^{\,a}_\beta(x),
  \qquad y>0,
  \label{eq:bcft-bulk-boundary-ope}
\end{equation}
where the boundary fields \(\widehat{\mathcal O}^{\,a}_\beta\) have
one-dimensional conformal weights \(h_\beta\).
\item Correlation functions on bordered Riemann surfaces with boundary
segments labelled by elements of \(\mathsf B\), satisfying locality,
reflection positivity in the unitary case, conformal covariance, and sewing
compatibility for gluing boundary and bulk circles.
\end{enumerate}
The same closed CFT can admit many inequivalent choices of \(\mathsf B\) and
sewing data.
\end{definition}

For a boundary equal to the real axis in the upper half-plane
\(\mathbb H=\{z=x+\ii y\mid y\ge0\}\), the stress tensor has components
\[
  T(z)=T_{zz}(z),\qquad \bar T(\bar z)=T_{\bar z\bar z}(\bar z).
\]
The component of energy flow through the boundary is proportional to
\(T(z)-\bar T(\bar z)\) at \(z=\bar z=x\).  A conformal boundary condition
sets this flux to zero:
\begin{equation}
  T(x)=\bar T(x),
  \qquad x\in\mathbb R,
  \label{eq:bcft-stress-tensor-gluing}
\end{equation}
as an identity inside boundary correlation functions away from boundary
operator insertions.

\paragraph{Preserved Virasoro algebra.}
Suppose a boundary condition on the upper half-plane obeys
\eqref{eq:bcft-stress-tensor-gluing}.  The conformal transformations
preserving the boundary are then generated by a single copy of the Virasoro
algebra.  In the closed-channel boundary-state convention the modes are
\begin{equation}
  L_n^{\partial}
  =
  L_n-\bar L_{-n}
  \label{eq:bcft-boundary-virasoro-generator}
\end{equation}
or equivalently by the diagonal combination of holomorphic and
antiholomorphic conformal vector fields on the half-plane.

Let \(v(z)\partial_z+\bar v(\bar z)\partial_{\bar z}\) be a conformal vector
field on the upper half-plane.  It preserves the real boundary precisely when
\(v(x)=\bar v(x)\) for real \(x\).  The Ward identity for an infinitesimal
transformation is the contour integral of
\[
  v(z)T(z)\dd z+\bar v(\bar z)\bar T(\bar z)\dd\bar z .
\]
Moving the contour to the boundary produces a term proportional to
\[
  v(x)\bigl(T(x)-\bar T(x)\bigr)\dd x,
\]
which vanishes by \eqref{eq:bcft-stress-tensor-gluing}.  On the cylinder,
the closed-channel boundary is a state at Euclidean time \(0\).  Reversing
the orientation of the antiholomorphic contour sends \(\bar L_n\) to
\(\bar L_{-n}\), so the annihilating generator is
\(L_n-\bar L_{-n}\).  These generators obey a single Virasoro algebra
because the two central terms cancel in the difference.

\section{Boundary States and Ishibashi States}

The same annulus can be quantized in two ways.  In the open channel, time
runs along the strip and the Hilbert space is \(\Hilb_{ab}\).  In the closed
channel, time runs across the strip and the two boundary components are
represented by states in the closed-string Hilbert space.  The equality of
these two descriptions is the annulus sewing constraint.

Let the closed Hilbert space of a diagonal rational CFT be
\begin{equation}
  \Hilb_{\mathrm{cl}}
  =
  \bigoplus_{i\in I}
  \mathcal H_i\otimes\overline{\mathcal H_i},
  \label{eq:diagonal-closed-hilbert-space-bcft}
\end{equation}
where \(\mathcal H_i\) are irreducible representations of the chiral algebra.
For an interval of length \(L\) and Euclidean time \(\beta\), define
\[
  q=\exp(-\pi\beta/L),
  \qquad
  \tilde q=\exp(-4\pi L/\beta).
\]
The same annulus amplitude must have the form
\begin{equation}
  Z_{ab}^{\mathrm{open}}(q)
  =
  \operatorname{Tr}_{\Hilb_{ab}}q^{L_0-c/24}
  =
  \langle a|
  \tilde q^{(L_0+\bar L_0-c/12)/2}
  |b\rangle
  =
  Z_{ab}^{\mathrm{closed}}(\tilde q).
  \label{eq:bcft-open-closed-annulus}
\end{equation}
The factor \(1/2\) in the closed Hamiltonian is a convention fixed by the
modular parameter of the annulus; changing cylinder coordinates changes both
sides by the same reparametrization.

\begin{definition}[Ishibashi state]
\label{def:ishibashi-state}
Assume \(\mathcal H_i\) is paired with its conjugate in
\eqref{eq:diagonal-closed-hilbert-space-bcft}.  Choose a homogeneous
orthonormal basis \(\{e_\alpha\}\) of \(\mathcal H_i\) and an antiunitary
map \(U_i:\mathcal H_i\to\overline{\mathcal H_i}\) intertwining the chiral
algebra with the antiholomorphic copy.  The convention is
\begin{equation}
  U_i L_m U_i^{-1}=\bar L_m
  \label{eq:ishibashi-antiunitary-convention}
\end{equation}
on the finite-energy domain.  For an additional chiral generator \(W^a_m\) of
spin \(s_a\), the preserved boundary condition is encoded by the corresponding
antiunitary adjoint convention
\begin{equation}
  U_i (W^a_m)^\dagger U_i^{-1}
  =
  (-1)^{s_a}\Omega(\bar W^a_{-m}),
  \label{eq:ishibashi-current-antiunitary-convention}
\end{equation}
where \(\Omega\) is the chosen automorphism of the chiral algebra.  The
Ishibashi state associated to \(i\) is the formal state
\begin{equation}
  |i\rangle\!\rangle
  =
  \sum_\alpha e_\alpha\otimes U_i e_\alpha .
  \label{eq:ishibashi-state-definition}
\end{equation}
It is not usually a normalizable vector in \(\Hilb_{\mathrm{cl}}\); it is a
distributional state whose overlaps are defined after inserting
\(\tilde q^{(L_0+\bar L_0-c/12)/2}\).
\end{definition}

\begin{proposition}[Ishibashi gluing]
\label{prop:ishibashi-gluing}
The state \(|i\rangle\!\rangle\) satisfies
\begin{equation}
  (L_n-\bar L_{-n})|i\rangle\!\rangle=0
  \label{eq:ishibashi-virasoro-gluing}
\end{equation}
for all \(n\in\mathbb Z\).  More generally, if \(W^a_m\) is a mode of a
chiral current of spin \(s_a\), the symmetry-preserving Ishibashi state
satisfies
\begin{equation}
  \left(W^a_m-(-1)^{s_a}\Omega(\bar W^a_{-m})\right)|i\rangle\!\rangle=0
  \label{eq:ishibashi-current-gluing}
\end{equation}
for the chosen automorphism \(\Omega\) of the chiral algebra.
\end{proposition}

\begin{proof}
Let
\[
  C_i=\sum_\alpha e_\alpha\otimes U_i e_\alpha
\]
denote the formal coevaluation vector.  We first record the finite-level
identity which controls the antiunitary bookkeeping.  If \(A\) is an operator
on the finite-energy domain of \(\mathcal H_i\) with adjoint \(A^\dagger\),
then
\begin{equation}
  (A\otimes 1)C_i=(1\otimes U_i A^\dagger U_i^{-1})C_i
  \label{eq:coevaluation-adjoint-identity}
\end{equation}
as a distributional vector, meaning after pairing against any
finite-level vector \(u\otimes U_iv\).  Indeed, using the Hermitian inner
product on \(\mathcal H_i\), written antilinear in the first argument, the
antiunitarity of \(U_i\), and the finite-level resolution of the identity
\(\sum_\alpha |e_\alpha\rangle\langle e_\alpha|=1\), the two pairings are
\[
  \sum_\alpha
  \langle u,Ae_\alpha\rangle\,\langle e_\alpha,v\rangle
  =
  \langle u,Av\rangle
  =
  \langle A^\dagger u,v\rangle
\]
and
\[
  \sum_\alpha
  \langle u,e_\alpha\rangle\,\langle A^\dagger e_\alpha,v\rangle
  =
  \langle A^\dagger u,v\rangle .
\]
Since finite-level vectors separate the graded dual completion in which
\(|i\rangle\!\rangle\) lives, \eqref{eq:coevaluation-adjoint-identity}
follows.  Taking \(A=L_n\) and using positive-energy unitarity
\(L_n^\dagger=L_{-n}\), together with
\eqref{eq:ishibashi-antiunitary-convention}, gives
\[
  (L_n\otimes 1)|i\rangle\!\rangle
  =
  (1\otimes \bar L_{-n})|i\rangle\!\rangle ,
\]
which is \eqref{eq:ishibashi-virasoro-gluing}.  For a spin-\(s_a\) current the
finite-level identity \eqref{eq:coevaluation-adjoint-identity}, applied to
\(A=W^a_m\), gives the adjoint mode in the reflected copy.  The boundary
reflection reverses the chiral coordinate and contributes the spin phase
\((-1)^{s_a}\); the chosen preserved diagonal current is then allowed to act
through the chiral-algebra automorphism \(\Omega\).  These three inputs give
\eqref{eq:ishibashi-current-gluing}.
\end{proof}

The overlap of two Ishibashi states is diagonal:
\begin{equation}
  \langle\!\langle i|
  \tilde q^{(L_0+\bar L_0-c/12)/2}
  |j\rangle\!\rangle
  =
  \delta_{ij}\chi_i(\tilde q),
  \label{eq:ishibashi-overlap-character}
\end{equation}
where \(\chi_i\) is the chiral character.  Thus a boundary state preserving
the chiral algebra has the expansion
\begin{equation}
  |a\rangle=\sum_i B_a{}^i\,|i\rangle\!\rangle
  \label{eq:bcft-boundary-state-expansion}
\end{equation}
with coefficients constrained by the open-channel interpretation.

\section{Cardy Consistency}

Assume the rational diagonal setting of
\eqref{eq:diagonal-closed-hilbert-space-bcft}.  The open-channel Hilbert
space between two elementary boundary conditions \(a,b\) must decompose as
\begin{equation}
  \Hilb_{ab}
  =
  \bigoplus_k n_{ab}{}^k\,\mathcal H_k,
  \qquad
  n_{ab}{}^k\in\mathbb Z_{\ge0}.
  \label{eq:open-channel-hilbert-space-decomposition}
\end{equation}
Using \(\chi_i(\tilde q)=\sum_k S_{ik}\chi_k(q)\), the annulus equality gives
\begin{equation}
  n_{ab}{}^k
  =
  \sum_i B_a{}^i B_b{}^i S_{ik}
  \label{eq:cardy-annulus-integrality-equation}
\end{equation}
for real unitary boundary coefficients.  Cardy consistency is the demand
that these numbers be nonnegative integers and that the resulting open
operator algebras sew consistently.

\paragraph{Cardy solution for the diagonal invariant.}
Let a unitary rational CFT have diagonal closed Hilbert space
\eqref{eq:diagonal-closed-hilbert-space-bcft}, unitary modular matrix \(S\),
and fusion coefficients \(N_{ab}{}^k\) obeying the Verlinde formula.  Then
the boundary states
\begin{equation}
  |a\rangle
  =
  \sum_i
  \frac{S_{ai}}{\sqrt{S_{0i}}}
  |i\rangle\!\rangle,
  \qquad a\in I,
  \label{eq:cardy-state-diagonal}
\end{equation}
satisfy the annulus integrality condition with
\begin{equation}
  n_{ab}{}^k=N_{ab}{}^k .
  \label{eq:cardy-open-spectrum-fusion}
\end{equation}

Substitute \(B_a{}^i=S_{ai}/\sqrt{S_{0i}}\) and
\(B_b{}^i=S_{bi}/\sqrt{S_{0i}}\) into
\eqref{eq:cardy-annulus-integrality-equation}.  Since the theory is unitary
and the Cardy labels are self-conjugate after replacing \(b\) by its
conjugate when necessary,
\[
  n_{ab}{}^k
  =
  \sum_i
  \frac{S_{ai}S_{bi}S_{ki}}{S_{0i}} .
\]
This is exactly the Verlinde formula for the fusion coefficient
\(N_{ab}{}^k\), with complex conjugation inserted on the label appropriate
to the orientation convention.  Fusion coefficients are dimensions of
three-point chiral block spaces, hence are nonnegative integers.

\section{Boundary OPE and Cardy--Lewellen Equations}
\label{sec:bcft-boundary-ope-sewing}

The annulus only detects the graded dimensions of open-channel Hilbert
spaces.  A BCFT also has boundary operator products.  If the boundary
condition is \(a\) to the left of an insertion and \(b\) to the right, write
\[
  \psi^{ab,\alpha}_i(x)
\]
for a boundary-condition-changing primary transforming in the chiral
representation \(i\).  The multiplicity label \(\alpha\) runs through a
basis of the corresponding space of boundary fields.

\begin{hypothesis}[Rational sewing hypotheses for this subsection]
\label{hyp:bcft-rational-sewing-hypotheses}
The explicit sewing equations in this subsection assume a unitary rational
chiral algebra with finitely many simple representations, semisimple
positive-energy modules, convergent genus-zero chiral conformal blocks,
nondegenerate boundary two-point pairings, and Moore--Seiberg fusing
isomorphisms satisfying the pentagon identity.  The diagonal Cardy example
also assumes the closed Hilbert space
\(\oplus_i\mathcal H_i\otimes\overline{\mathcal H_i}\).  Outside these
hypotheses the formulas remain useful local models, but they are not a proof
of analytic sewing.
\end{hypothesis}

Under Hypothesis~\ref{hyp:bcft-rational-sewing-hypotheses}, for ordered
boundary conditions \(a,b,c\), the boundary OPE has the form
\begin{equation}
  \psi^{ab,\alpha}_i(x)\psi^{bc,\beta}_j(0)
  \sim
  \sum_{k,\gamma}
  C\!\begin{bmatrix} a&b&c\\ i&j&k \end{bmatrix}_{\alpha\beta}^{\gamma}
  x^{h_k-h_i-h_j}
  \psi^{ac,\gamma}_k(0)+\cdots .
  \label{eq:boundary-changing-ope}
\end{equation}
The coefficient is not just a number attached to three labels; it depends on
normalizations of boundary two-point functions and on a choice of basis in
each chiral block space.  A change of those bases conjugates the coefficient
tensors, while the sewing equations below are invariant statements.

For four boundary fields on a disk, with boundary conditions
\[
  a\longrightarrow b\longrightarrow c\longrightarrow d
  \longrightarrow a ,
\]
there are two ways to perform the OPE.  In a rational theory the two chiral
decompositions are related by the fusing matrix \(F\).  The resulting
Cardy--Lewellen boundary sewing equation, with dual-basis pairings of chiral
blocks suppressed from the notation, is
\begin{align}
  &\sum_{p,\mu}
  C\!\begin{bmatrix} a&b&c\\ i&j&p \end{bmatrix}_{\alpha\beta}^{\mu}
  C\!\begin{bmatrix} a&c&d\\ p&k&\ell \end{bmatrix}_{\mu\gamma}^{\delta}
  F_{pq}\!\begin{bmatrix} i&j\\ k&\ell \end{bmatrix}
  \notag\\
  &\qquad =
  \sum_{\eta}
  C\!\begin{bmatrix} b&c&d\\ j&k&q \end{bmatrix}_{\beta\gamma}^{\eta}
  C\!\begin{bmatrix} a&b&d\\ i&q&\ell \end{bmatrix}_{\alpha\eta}^{\delta}.
  \label{eq:cardy-lewellen-boundary-sewing}
\end{align}
The displayed formula suppresses the standard dual-basis pairings of chiral
blocks.  Its invariant content is that boundary OPE coefficients form an
associative multiplication object in the chiral tensor category, with the
associator supplied by the Moore--Seiberg fusing isomorphism.

\paragraph{Cardy boundary fields from fusion.}
In the diagonal Cardy solution, the multiplicity of boundary primaries from
\(a\) to \(b\) is
\begin{equation}
  \dim\operatorname{span}\{\psi^{ab}_i\}=N_{ia}{}^b .
  \label{eq:cardy-boundary-field-multiplicity}
\end{equation}
If the boundary three-point coefficients are chosen to be the normalized
fusing matrices of the chiral theory, then
\eqref{eq:cardy-lewellen-boundary-sewing} is precisely the pentagon identity
for the fusing matrices.

The open spectrum for Cardy boundary labels is
\(\Hilb_{ab}\simeq\oplus_iN_{ab}{}^i\mathcal H_i\).  After conjugating
orientation conventions, this is equivalently
\eqref{eq:cardy-boundary-field-multiplicity}.  Boundary three-point
functions are chiral three-point blocks on the doubled disk.  The operation
of composing two boundary OPEs is therefore the operation of changing
parenthesization in a four-point chiral block.  The Moore--Seiberg fusing
matrix is exactly the transition matrix between the two parenthesizations.
The pentagon identity says that all ways of changing parenthesization in a
five-point block agree.  Specializing the pentagon to the four-boundary-field
factorization diagram gives \eqref{eq:cardy-lewellen-boundary-sewing}.

The disk one-point coefficient of a spinless bulk primary is another piece
of sewing data.  In the diagonal Cardy case the coefficient of the identity
boundary field is proportional to
\[
  B_a{}^i=\frac{S_{ai}}{\sqrt{S_{0i}}}.
\]
The Verlinde eigenvalue
\begin{equation}
  \lambda_a(i)=\frac{S_{ia}}{S_{0a}}
  \label{eq:cardy-fusion-ring-character}
\end{equation}
is a one-dimensional character of the fusion ring:
\begin{equation}
  \lambda_a(i)\lambda_a(j)
  =
  \sum_k N_{ij}{}^k\,\lambda_a(k).
  \label{eq:fusion-ring-character-identity}
\end{equation}
Indeed, the Verlinde formula says that the matrices \(N_i\) are diagonalized
by \(S\), with eigenvalues \(S_{ia}/S_{0a}\).  The Cardy disk coefficient is
the same eigenvalue with the normalization required by the bulk two-point
function:
\[
  B_a{}^i
  =
  \frac{S_{0a}}{\sqrt{S_{0i}}}\,\lambda_a(i).
\]
Thus the disk one-point functions solve the genus-zero bulk-boundary sewing
equation only after this normalization has been separated from the physical
two-point convention.  The useful coordinate for the disk sewing equation is
\begin{equation}
  \widehat B_a(i)
  :=
  \frac{\sqrt{S_{0i}}}{S_{0a}}\,B_a{}^i
  =
  \lambda_a(i).
  \label{eq:cardy-normalized-classifying-coordinate}
\end{equation}
In the diagonal Cardy theory the bulk fields which can appear in the boundary
identity channel form the fusion algebra object with basis \(\{e_i\}\) and
product
\begin{equation}
  e_i e_j=\sum_k N_{ij}{}^k e_k .
  \label{eq:cardy-classifying-algebra-product}
\end{equation}
The boundary condition \(a\) defines the linear functional
\(\chi_a(e_i)=\widehat B_a(i)\).  The disk sewing constraint obtained by
comparing the two orders of operation, first fusing two bulk fields and then
taking the boundary identity channel or first taking the two identity
bulk-boundary channels and multiplying the results, is the multiplicativity
condition
\begin{equation}
  \chi_a(e_i e_j)
  =
  \chi_a(e_i)\chi_a(e_j).
  \label{eq:cardy-disk-two-bulk-classifying-sewing}
\end{equation}
Substituting \eqref{eq:cardy-classifying-algebra-product} and
\(\chi_a(e_i)=\lambda_a(i)\) gives precisely
\eqref{eq:fusion-ring-character-identity}.  This is a statement about the
normalized classifying coordinate.  It is not the assertion that the raw
bulk OPE constants are fusion coefficients, nor that all metric factors have
disappeared; those factors have been absorbed into the chosen basis of the
classifying algebra and into the normalization of the identity boundary
field.  In a non-diagonal rational CFT the same low-genus sewing constraint
is formulated in the classifying algebra associated to the chosen Frobenius
algebra object; in a nonrational theory there need not be a
finite-dimensional classifying algebra at all, and the finite character
calculation must be replaced by an analytic statement about the relevant
continuous spectrum.

\section{Bulk-Boundary Normalization and Direct Sums}
\label{sec:bcft-bulk-boundary-normalization-direct-sums}

The boundary state coefficient, the disk one-point function, and the
identity term in the bulk-boundary OPE are often denoted by the same symbol.
That is harmless only after the two-point metric and the identity
normalization have been fixed.  We record the convention because otherwise
Cardy formulas can appear to disagree by factors of \(S_{0i}\).

Assume that bulk spinless primaries \(\phi_i\) are normalized on the plane by
\begin{equation}
  \langle \phi_i(z,\bar z)\phi_j(0,0)\rangle
  =
  \frac{D_{ij}}{|z|^{4h_i}},
  \qquad
  D_{ij}=0\ \text{unless}\ h_i=h_j,
  \label{eq:bcft-bulk-two-point-metric}
\end{equation}
with inverse matrix \(D^{ij}\).  For a boundary condition \(a\), write the
upper-half-plane one-point function as
\begin{equation}
  \langle \phi_i(x+\ii y,x-\ii y)\rangle_a
  =
  \frac{U^a_i}{(2y)^{2h_i}},
  \qquad y>0 .
  \label{eq:bcft-disk-one-point-lowered}
\end{equation}
If the boundary identity \(\mathbf 1_a\) has expectation value \(1\), the
identity term in the bulk-boundary OPE is
\begin{equation}
  \phi_i(x+\ii y,x-\ii y)
  =
  R^a_{i0}(2y)^{-2h_i}\mathbf 1_a+\cdots ,
  \qquad
  R^a_{i0}=U^a_i .
  \label{eq:bcft-bulk-boundary-identity-coefficient}
\end{equation}
Equivalently, if one raises the bulk label by defining
\[
  U_a{}^j=D^{ji}U^a_i ,
\]
then the same coefficient is
\begin{equation}
  R^a_{i0}=D_{ij}U_a{}^j .
  \label{eq:bcft-r-u-one-point-relation}
\end{equation}
Thus \eqref{eq:cardy-state-diagonal} determines a physical one-point
function only after the Ishibashi overlap normalization and the bulk
two-point metric have both been declared.

\begin{remark}[Direct sums and Chan--Paton matrix units]
\label{rem:bcft-direct-sum-chan-paton}
Let \(a,b\) be boundary conditions in a unitary BCFT, and suppose finite
direct sums of boundary conditions are allowed.  For
\[
  A=a^{\oplus n},\qquad B=b^{\oplus m},
\]
with no additional boundary fields beyond the direct-sum completion, the
open Hilbert space and boundary fields are
\begin{equation}
  \Hilb_{AB}
  \simeq
  \Hilb_{ab}\otimes \operatorname{Mat}_{n\times m}(\mathbb C).
  \label{eq:bcft-chan-paton-hilbert-space}
\end{equation}
If \(E_{rs}\) denotes the standard matrix unit, boundary OPEs obey
\begin{equation}
  (\psi_\alpha\otimes E_{rs})(x)
  (\psi_\beta\otimes E_{tu})(0)
  =
  \delta_{st}
  \bigl(\psi_\alpha(x)\psi_\beta(0)\bigr)\otimes E_{ru}.
  \label{eq:bcft-chan-paton-ope}
\end{equation}
Consequently
\begin{equation}
  Z_{A B}^{\mathrm{open}}(q)=nm\,Z_{ab}^{\mathrm{open}}(q),
  \qquad
  g_{a^{\oplus n}}=n\,g_a .
  \label{eq:bcft-chan-paton-annulus-g}
\end{equation}
Indeed, a direct sum boundary condition is a finite biproduct in the boundary
category.  An interval with left boundary \(a^{\oplus n}\) and right boundary
\(b^{\oplus m}\) therefore has one copy of \(\Hilb_{ab}\) for each ordered pair
of summands.  Choosing the standard bases of the multiplicity spaces identifies
this direct sum with \eqref{eq:bcft-chan-paton-hilbert-space}.  Locality on the
boundary forces composition to match endpoint labels: an operator carrying the
right endpoint index \(s\) can compose with a second operator whose left
endpoint is \(t\) only when \(s=t\).  This is precisely the multiplication rule
for matrix units, giving \eqref{eq:bcft-chan-paton-ope}.  Taking the graded
trace over \(\Hilb_{ab}\otimes\operatorname{Mat}_{n\times m}(\mathbb C)\) gives
the factor \(nm\) in the open annulus.  In the closed channel, the boundary
state of \(a^{\oplus n}\) is the sum of \(n\) identical boundary states, so its
vacuum overlap is \(n g_a\).
\end{remark}

\section{Compact Free Boson Boundaries}

Let \(X\sim X+2\pi R\) be a compact free boson with
\[
  X_L(z)X_L(0)\sim-\log z,
  \qquad
  X_R(\bar z)X_R(0)\sim-\log\bar z .
\]
Circle states are labelled by momentum \(m\in\mathbb Z\) and winding
\(w\in\mathbb Z\), with
\begin{equation}
  k_L=\frac mR+wR,
  \qquad
  k_R=\frac mR-wR,
  \qquad
  h=\frac12k_L^2+N,\quad
  \bar h=\frac12k_R^2+\bar N .
  \label{eq:compact-boson-momenta-bcft}
\end{equation}
The oscillator modes obey
\([a_m,a_n]=m\delta_{m+n,0}\) and similarly for \(\bar a_n\).

On the upper half-plane, Neumann and Dirichlet boundary conditions are
\[
  \partial_yX|_{y=0}=0,
  \qquad
  X|_{y=0}=x_0 .
\]
In closed-channel boundary-state language they become
\begin{align}
  (a_n+\bar a_{-n})|N,\theta\rangle&=0,
  \label{eq:neumann-oscillator-gluing}\\
  (a_n-\bar a_{-n})|D,x_0\rangle&=0,
  \label{eq:dirichlet-oscillator-gluing}
\end{align}
for all \(n\in\mathbb Z\).  The zero-mode parts are the important point:
\[
  \text{Neumann: } k_L+k_R=0\iff m=0,
  \qquad
  \text{Dirichlet: } k_L-k_R=0\iff w=0 .
\]
Thus Neumann boundary states sum over winding and carry a Wilson-line phase
\(\theta\), while Dirichlet boundary states sum over momentum and carry a
position phase \(x_0\):
\begin{align}
  |N,\theta\rangle
  &=
  \mathcal N_N
  \sum_{w\in\mathbb Z}
  e^{\ii w\theta}
  \exp\!\left(
    -\sum_{n\ge1}\frac1n a_{-n}\bar a_{-n}
  \right)|0,w\rangle,
  \label{eq:compact-boson-neumann-state}\\
  |D,x_0\rangle
  &=
  \mathcal N_D
  \sum_{m\in\mathbb Z}
  e^{\ii m x_0/R}
  \exp\!\left(
    \sum_{n\ge1}\frac1n a_{-n}\bar a_{-n}
  \right)|m,0\rangle .
  \label{eq:compact-boson-dirichlet-state}
\end{align}
The constants \(\mathcal N_N,\mathcal N_D\) are fixed by the annulus
normalization.  With the oscillator trace normalized by \(1/\eta(q)\), one
may choose \(\mathcal N_N=\sqrt R\) and
\(\mathcal N_D=1/\sqrt{2R}\) in the present momentum convention; changing
the normalization of the circle volume rescales both boundary states and
the open-string ground-state density by the same factor.

The radius-duality map
\[
  R\longmapsto R^{-1},
  \qquad
  m\longleftrightarrow w,
  \qquad
  X_R\longmapsto -X_R
\]
exchanges the gluing equations
\eqref{eq:neumann-oscillator-gluing} and
\eqref{eq:dirichlet-oscillator-gluing}.  It sends a Neumann brane with
Wilson line \(\theta\) to a Dirichlet brane at the dual position \(\theta\),
and sends a Dirichlet position \(x_0\) to a Wilson line on the dual circle.
The transformation \(X_R\mapsto -X_R\) sends
\(\bar a_n\mapsto-\bar a_n\), while leaving \(a_n\) unchanged.  Therefore
\(a_n+\bar a_{-n}\) becomes \(a_n-\bar a_{-n}\).  On zero modes,
\((m,w)\mapsto(w,m)\) exchanges the condition \(m=0\) with \(w=0\).  The
phases in \eqref{eq:compact-boson-neumann-state} and
\eqref{eq:compact-boson-dirichlet-state} then exchange Wilson-line and
position data.

\section{Majorana Fermion and Ising Boundary States}

The chiral Majorana field has modes
\[
  \psi(z)=\sum_r\psi_r z^{-r-1/2},
  \qquad
  \{\psi_r,\psi_s\}=\delta_{r+s,0}.
\]
In the Neveu--Schwarz sector \(r\in\mathbb Z+1/2\), while in the Ramond
sector \(r\in\mathbb Z\).  A conformal fermion boundary condition is
\begin{equation}
  \left(\psi_r+\ii\eta\,\bar\psi_{-r}\right)|B,\eta\rangle=0,
  \qquad
  \eta=\pm1 .
  \label{eq:majorana-fermion-gluing}
\end{equation}
The stress-tensor gluing follows because the fermion stress tensor is
quadratic in \(\psi\), but the converse is weaker: a Virasoro boundary
condition need not preserve the fermion field itself.

For the bosonic Ising model, the three Virasoro primaries are
\[
  \mathbf 1,\qquad \epsilon,\qquad \sigma
\]
with weights \(0,1/2,1/16\).  In the ordered basis
\((\mathbf 1,\epsilon,\sigma)\), the modular matrix is
\begin{equation}
  S_{\mathrm{Ising}}
  =
  \frac12
  \begin{pmatrix}
    1&1&\sqrt2\\
    1&1&-\sqrt2\\
    \sqrt2&-\sqrt2&0
  \end{pmatrix}.
  \label{eq:ising-modular-s-matrix-bcft}
\end{equation}
Cardy's formula gives the three elementary boundary states
\begin{align}
  |+\rangle
  &=
  \frac1{\sqrt2}|\mathbf 1\rangle\!\rangle
  +
  \frac1{\sqrt2}|\epsilon\rangle\!\rangle
  +
  2^{-1/4}|\sigma\rangle\!\rangle,
  \label{eq:ising-plus-boundary-state}\\
  |-\rangle
  &=
  \frac1{\sqrt2}|\mathbf 1\rangle\!\rangle
  +
  \frac1{\sqrt2}|\epsilon\rangle\!\rangle
  -
  2^{-1/4}|\sigma\rangle\!\rangle,
  \label{eq:ising-minus-boundary-state}\\
  |f\rangle
  &=
  |\mathbf 1\rangle\!\rangle
  -
  |\epsilon\rangle\!\rangle .
  \label{eq:ising-free-boundary-state}
\end{align}
The labels \(+\), \(-\), and \(f\) denote the two fixed spin boundaries and
the free boundary.  The annulus spectra are the Ising fusion rules:
\[
  \Hilb_{++}\simeq\mathcal H_{\mathbf 1},\qquad
  \Hilb_{+-}\simeq\mathcal H_{\epsilon},\qquad
  \Hilb_{ff}\simeq\mathcal H_{\mathbf 1}\oplus\mathcal H_{\epsilon},
\]
and
\[
  \Hilb_{+f}\simeq\mathcal H_{\sigma},
  \qquad
  \Hilb_{-f}\simeq\mathcal H_{\sigma}.
\]
The NS/R bookkeeping is visible in the fermionic presentation: the
\(\mathbf 1\) and \(\epsilon\) Virasoro modules come from NS fermion
parity sectors, while \(\sigma\) is the Ramond spin field.  A spin CFT keeps
the spin-structure dependence explicitly; the bosonic Ising CFT is obtained
after the corresponding fermion-parity gauging.

\paragraph{Ising boundary-changing OPE constants.}
Identify the Cardy labels
\[
  +=\mathbf 1,\qquad -=\epsilon,\qquad f=\sigma .
\]
Let \(\psi_i^{ab}\) denote the boundary primary of chiral type \(i\) between
boundary labels \(a\) and \(b\), when the fusion multiplicity
\(N_{ia}{}^b\) is one.  Use the Cardy boundary-field basis determined by the
Ising fusing symbols, and fix the residual scalar normalizations of
one-dimensional boundary-field spaces so that the one-channel products
displayed below have coefficient \(1\).  The only products below with two
allowed final channels are the products in which two \(\sigma\)-type
boundary-changing fields leave and return to the free boundary through a
fixed boundary.  For \(x>0\),
\begin{align}
  \psi_\sigma^{f+}(x)\psi_\sigma^{+f}(0)
  &\sim
  2^{-1/2}x^{-1/8}\mathbf 1_f
  +
  2^{-1/2}x^{3/8}\psi_\epsilon^{ff}(0)+\cdots ,
  \label{eq:ising-bcc-plus-through-free}\\
  \psi_\sigma^{f-}(x)\psi_\sigma^{-f}(0)
  &\sim
  2^{-1/2}x^{-1/8}\mathbf 1_f
  -
  2^{-1/2}x^{3/8}\psi_\epsilon^{ff}(0)+\cdots .
  \label{eq:ising-bcc-minus-through-free}
\end{align}
The opposite orientation has no channel multiplicity:
\begin{align}
  \psi_\sigma^{+f}(x)\psi_\sigma^{f+}(0)
  &\sim x^{-1/8}\mathbf 1_+ +\cdots,\\
  \psi_\sigma^{-f}(x)\psi_\sigma^{f-}(0)
  &\sim x^{-1/8}\mathbf 1_- +\cdots,\\
  \psi_\sigma^{+f}(x)\psi_\sigma^{f-}(0)
  &\sim x^{3/8}\psi_\epsilon^{+-}(0)+\cdots,\\
  \psi_\sigma^{-f}(x)\psi_\sigma^{f+}(0)
  &\sim x^{3/8}\psi_\epsilon^{-+}(0)+\cdots .
\end{align}
The multiplicities follow from
\eqref{eq:cardy-boundary-field-multiplicity} and the Ising fusion rules.
Thus the nontrivial boundary-changing fields are
\[
  \psi_\epsilon^{+-},\quad \psi_\epsilon^{-+},\quad
  \psi_\epsilon^{ff},\quad
  \psi_\sigma^{+f},\quad \psi_\sigma^{f+},\quad
  \psi_\sigma^{-f},\quad \psi_\sigma^{f-}.
\]
In the diagonal Cardy case the boundary OPE constants are the fusing
coefficients of the chiral theory in the chosen basis.  The only
nontrivial Ising fusing matrix needed here is the recoupling of three
\(\sigma\) lines with total charge \(\sigma\):
\begin{equation}
  F^{\sigma\sigma\sigma}_{\sigma}
  =
  \frac1{\sqrt2}
  \begin{pmatrix}
    1&1\\
    1&-1
  \end{pmatrix},
  \label{eq:ising-f-symbol-sigma-sigma-sigma}
\end{equation}
where the row records the intermediate fixed boundary
\(\mathbf 1\) or \(\epsilon\), and the column records the final channel
\(\mathbf 1_f\) or \(\psi_\epsilon^{ff}\) on the free boundary.  Equations
\eqref{eq:ising-bcc-plus-through-free}--\eqref{eq:ising-bcc-minus-through-free}
are precisely the two rows of
\eqref{eq:ising-f-symbol-sigma-sigma-sigma}.  The powers of \(x\) are the
boundary OPE powers \(h_k-h_i-h_j\):
\[
  0-\frac1{16}-\frac1{16}=-\frac18,\qquad
  \frac12-\frac1{16}-\frac1{16}=\frac38 .
\]
All remaining displayed OPEs have one-dimensional fusion spaces.  Their
coefficients have been fixed to be \(1\) by the residual scalar
normalizations specified in the statement; changing those normalizations
rescales the corresponding raw three-point constants but not the sewing
class represented by \eqref{eq:ising-f-symbol-sigma-sigma-sigma}.
Thus rescaling the fields \(\psi_i^{ab}\) changes the displayed constants by
the corresponding two-point normalization factors.  The invariant content is
the Ising fusing matrix and, in particular, the relative sign of the
\(\psi_\epsilon^{ff}\) coefficient for the two fixed boundaries.

\paragraph{A finite Ising sewing cell.}
The preceding constants give a completely explicit instance of the
Cardy--Lewellen equation.  Let \(r,s\in\{+,-\}\), and consider four boundary
fields arranged around the disk as
\[
  f\xrightarrow{\ \psi_\sigma^{fr}\ }r
  \xrightarrow{\ \psi_\sigma^{rf}\ }f
  \xrightarrow{\ \psi_\sigma^{fs}\ }s
  \xrightarrow{\ \psi_\sigma^{sf}\ }f .
\]
Write
\[
  \psi_\sigma^{fr}(x)\psi_\sigma^{rf}(0)
  \sim
  v_r(\mathbf 1)\,x^{-1/8}\mathbf 1_f
  +
  v_r(\epsilon)\,x^{3/8}\psi_\epsilon^{ff}(0)+\cdots .
\]
Equations~\eqref{eq:ising-bcc-plus-through-free} and
\eqref{eq:ising-bcc-minus-through-free} give
\[
  v_+=2^{-1/2}(1,1),\qquad
  v_-=2^{-1/2}(1,-1),
\]
where the two entries are the \(\mathbf 1\) and \(\epsilon\) channels on the
free boundary.  If the two adjacent pairs are fused first through the free
boundary, the coefficient multiplying the common free-boundary intermediate
channel is
\[
  \sum_{p\in\{\mathbf 1,\epsilon\}} v_r(p)v_s(p).
\]
If instead one first fuses across the fixed-boundary side, the corresponding
open-sector morphism space is \(\operatorname{Hom}(U_r,U_s)\), which has
dimension \(\delta_{rs}\).  The sewing equality for this cell is therefore
\begin{equation}
  \sum_{p\in\{\mathbf 1,\epsilon\}} v_r(p)v_s(p)=\delta_{rs}.
  \label{eq:ising-boundary-four-point-sewing-cell}
\end{equation}
This is the matrix identity
\[
  F_{\sigma}^{\sigma\sigma\sigma}
  \left(F_{\sigma}^{\sigma\sigma\sigma}\right)^{T}
  =
  \mathbf 1
\]
for the nontrivial Ising fusing matrix
\eqref{eq:ising-f-symbol-sigma-sigma-sigma}.  In this example, the
abstract module associativity statement in
\eqref{eq:cardy-lewellen-boundary-sewing} reduces to a visible orthogonality
relation between the two fixed-boundary rows.  The calculation does not prove
analytic sewing for a general BCFT; it identifies exactly which finite
algebraic identity the sewing theorem uses in the diagonal Ising cell.

\section{Nonrational Example: Liouville Boundary States}
\label{sec:bcft-liouville-boundary-states}

Chapter~\ref{ch:liouville-cft}, Section~\ref{sec:liouville-boundary-states},
constructs the Liouville boundary states used as the basic nonrational test
case for this chapter.  The closed state space is no longer a finite
diagonal sum; in the scattering normalization used here it is a direct integral
\[
  \Hilb_{\mathrm{cl}}^{\mathrm{L}}
  =
  \int_0^\infty
  \mathcal V_{Q/2+\ii P}\otimes
  \overline{\mathcal V}_{Q/2+\ii P}\,\dd P ,
\]
and the Ishibashi label \(i\) is replaced by the continuous momentum
\(P\).  A boundary state is therefore a distributional wavefunction
\(\Psi_a(P)\) multiplying \(|P\rangle\!\rangle\), not a finite vector of
Cardy coefficients.

For the FZZT family, the wavefunction is
\[
  \Psi_s(P)
  =
  \frac1{\ii P}
  \left(\pi\mu\,\gamma(b^2)\right)^{-\ii P/b}
  \Gamma(1+2\ii bP)\Gamma(1+2\ii P/b)
  \cos(2\pi sP),
\]
where \(\gamma(x)=\Gamma(x)/\Gamma(1-x)\), up to the common normalization and
measure convention fixed in
\eqref{eq:fzzt-boundary-state}.  The discrete ZZ states are finite
differences at imaginary FZZT parameters, so their wavefunctions contain
\[
  2\sinh(2\pi mP/b)\sinh(2\pi nbP).
\]
This is the nonrational replacement for the finite modular \(S\)-matrix
entry in Cardy's rational formula: annulus consistency is expressed by
integral kernels and hyperbolic identities, followed by an analytic
statement about the resulting spectral measure.

\paragraph{Hyperbolic kernels behind Liouville annuli.}
For formal variables \(X,Y\) and for every \(n\in\mathbb Z_{>0}\),
\begin{align}
  \cosh(X+Y)-\cosh(X-Y)
  &=
  2\sinh X\sinh Y,
  \label{eq:liouville-zz-hyperbolic-identity}
  \\
  \frac{\sinh(nX)}{\sinh X}
  &=
  \sum_{j=0}^{n-1}\exp\!\left((n-1-2j)X\right).
  \label{eq:liouville-degenerate-shift-sum}
\end{align}
The first identity is the exact algebra turning an imaginary FZZT difference
into a ZZ wavefunction.  The second identity is the exact algebra turning
degenerate Liouville annulus kernels into a finite sum of shifted
nondegenerate characters.  The verification is elementary algebra:
\(\cosh Z=(\ee^Z+\ee^{-Z})/2\), and multiplying the finite sum in
\eqref{eq:liouville-degenerate-shift-sum} by
\(2\sinh X=\ee^X-\ee^{-X}\) leaves only the telescoping endpoints
\(\ee^{nX}-\ee^{-nX}\).  The identities are deliberately only the algebraic
part of the story.  A complete Liouville BCFT construction must still specify
the direct-integral topology, the domains of boundary-condition-changing
fields, the treatment of poles encountered under contour motion, and the
positivity or distribution interpretation of the annulus spectral density.
Those are analytic sewing requirements, not consequences of the finite
hyperbolic identities alone.

\paragraph{Analytic sewing datum for nonrational boundaries.}
The preceding wavefunctions become a boundary conformal field theory only
after they are placed in a functional-analytic sewing datum.  One convenient
way to state the target is the following.  First choose, inside the closed
Liouville direct integral, a dense nuclear test space
\(\mathcal S_{\rm cl}\) whose vectors have rapidly decreasing smooth
dependence on \(P\) and lie levelwise in the smooth Virasoro-module
vectors.  Boundary states in this setting are continuous linear functionals
\[
  |a\rangle\in\mathcal S_{\rm cl}' .
\]
The FZZT and ZZ formulas above specify candidate distributional coordinates
of such functionals.  The regulated closed-channel annulus is then the
distributional pairing
\begin{equation}
  Z_{ab}^{\rm cl}(\tilde q)
  =
  \bigl\langle a\bigm|
  \tilde q^{(L_0+\bar L_0-c/12)/2}
  \bigm|b\bigr\rangle ,
  \qquad 0<|\tilde q|<1,
  \label{eq:liouville-bcft-closed-annulus-pairing}
\end{equation}
initially defined by testing the two distributions against the trace-class
heat kernel on each positive-energy Virasoro module, with the continuous
\(P\)-integration still treated distributionally, and then by analytic
continuation in the boundary parameters when poles are crossed.

Second, for every ordered pair of boundary conditions \(a,b\), one must
construct an open-channel positive-energy Hilbert space \(\Hilb_{ab}\) and a
dense common smooth domain \(\mathcal D_{ab}\subset\Hilb_{ab}\).  Annulus
positivity is the existence of positive spectral measures \(\nu_{ab}\) and
multiplicity fields \(\mathcal M_{ab}(h)\) such that
\begin{equation}
  Z_{ab}^{\rm op}(q)
  =
  \int_{\mathbb R_{\ge0}}
  \operatorname{Tr}_{\mathcal M_{ab}(h)}
  q^{h-c/24}\,\dd\nu_{ab}(h)
  \label{eq:liouville-bcft-open-annulus-measure}
\end{equation}
agrees with \eqref{eq:liouville-bcft-closed-annulus-pairing} in the common
domain of convergence, with the usual matrix-positivity condition for finite
linear combinations of boundary labels.  In degenerate cases the measure may
have a discrete part.  The ZZ finite differences are meaningful precisely
because the contour prescription turns the corresponding closed-channel
functional into such a positive open-channel spectral datum; the hyperbolic
telescoping identity \eqref{eq:liouville-degenerate-shift-sum} is its finite
algebraic shadow.

Third, boundary-condition-changing fields must be defined as
operator-valued distributions on these domains.  A primary changing
\(a\) to \(b\), inserted on a boundary segment, should give multilinear
matrix elements
\[
  \bigl\langle v_0,\,
  \psi^{ab}_{\lambda_1}(x_1)
  \cdots
  \psi^{cd}_{\lambda_n}(x_n)\,
  v_n\bigr\rangle ,
  \qquad x_1>\cdots>x_n ,
\]
with \(v_0,v_n\) in the appropriate open-channel smooth domains or their
duals.  The labels \(\lambda_i\) may be continuous, discrete, or resonant.
The OPE coefficients are kernels between spectral measures, so their
associativity is an equality of distributions after smearing in the spectral
variables.  This is the nonrational replacement for finite sums over
boundary primaries in \eqref{eq:boundary-changing-ope}.

Fourth, the analytic continuation in FZZT parameters and the passage to ZZ
states must include a contour and residue prescription.  When the boundary
parameter is moved to an imaginary degenerate value, the \(P\)-contour in
closed-channel pairings can encounter poles of the wavefunctions and of
bulk-boundary or boundary-boundary structure kernels.  A sewing datum must
state whether the amplitude is defined by contour deformation, by principal
value plus residues, or by an equivalent distributional boundary value, and
must verify that the same prescription is used in all channels.  A
channel-dependent prescription defines incompatible disk or annulus
amplitudes even when each individual channel satisfies the displayed
finite-difference algebra.

Finally, bordered-surface sewing requires continuous multilinear maps on the
chosen test spaces.  Closed-circle gluing is integration over
\[
  \int_0^\infty \dd P\,\mu_{\rm L}(P)
\]
with the Liouville Plancherel measure and descendant two-point pairing;
open-interval gluing is integration or summation over the open spectral
measure \(\nu_{ab}\).  In a plumbing chart the resulting series or integral
must converge for \(|q_e|<1\) and must be compatible with changing the order
of gluing.  Anomaly determinant lines, if present, must be tensored into
these sewing maps.  This paragraph is therefore a definition of the analytic
burden carried by a nonrational BCFT construction.  The Liouville FZZT/ZZ
formulas supply strong candidate coordinates and many exact residue
identities; proving the full bordered-surface functor requires the analytic
data above.

\section{Sewing, Boundary Entropy, and Defect Interfaces}

Section~\ref{sec:bcft-boundary-ope-sewing} displayed the boundary
four-point sewing equation and the fusion-ring character equation for disk
one-point functions in the diagonal Cardy case.  A full BCFT also has
bulk-boundary OPE coefficients with nontrivial boundary fields, bulk
three-point coefficients, and compatibility with defect endpoints.  These
data obey additional compatibility equations from sewing a disk with two
bulk fields and one boundary field in different orders, and from comparing
bulk, boundary, and defect factorizations.

\begin{quotedtheorem}[Rational Cardy--Lewellen construction boundary]
\label{qthm:cardy-lewellen-sewing-status}
For rational chiral data satisfying the analytic sewing hypotheses of a
modular functor, a special symmetric Frobenius algebra object \(A\) in the
chiral modular tensor category produces a full rational CFT with boundaries;
boundary conditions are left \(A\)-modules, and the genus-zero
Cardy--Lewellen equations follow from the Frobenius and module identities.
Converses and classification statements require additional equivalence data
and are not used below.  This chapter uses only the local finite-categorical
consequences developed in the following paragraphs, not a proof of
nonrational or continuous-spectrum BCFT sewing.
\end{quotedtheorem}

\paragraph{Algebraic core of the rational boundary construction.}
The quoted theorem has a concrete finite-categorical mechanism.  Let
\(\mathsf C\) be the semisimple ribbon tensor category of chiral
representations, with tensor product \(\otimes\), unit object
\(\mathbf 1\), associator \(\alpha\), braiding \(c\), and representatives
\(\{U_i\}_{i\in I}\) of simple objects.  A rational full CFT with a chosen
bulk modular invariant and a chosen family of topological boundary
conditions is encoded algebraically by an object \(A\in\mathsf C\) equipped
with maps
\[
  m:A\otimes A\to A,\qquad
  \eta:\mathbf 1\to A,\qquad
  \Delta:A\to A\otimes A,\qquad
  \varepsilon:A\to\mathbf 1 .
\]
The multiplication and unit make \(A\) an algebra object:
\[
  m(m\otimes\operatorname{id}_A)
  =
  m(\operatorname{id}_A\otimes m)\alpha_{A,A,A},
  \qquad
  m(\eta\otimes\operatorname{id}_A)=\operatorname{id}_A
  =
  m(\operatorname{id}_A\otimes\eta).
\]
The comultiplication and counit obey the dual coalgebra identities, and the
Frobenius compatibility is
\begin{equation}
  (m\otimes\operatorname{id}_A)
  (\operatorname{id}_A\otimes\Delta)
  =
  \Delta m
  =
  (\operatorname{id}_A\otimes m)(\Delta\otimes\operatorname{id}_A),
  \label{eq:bcft-frobenius-identity-object}
\end{equation}
with associators inserted in the unique type-correct positions.  The word
special means that \(m\Delta\) is a nonzero scalar multiple of
\(\operatorname{id}_A\), so cutting out a short \(A\)-ribbon segment does not
change amplitudes except for a fixed normalization.  The word symmetric
means that the two natural maps \(A\to A^\vee\) obtained from
\(\varepsilon m\) agree; this is the categorical form of a nondegenerate
two-point pairing compatible with orientation reversal.  For a local
oriented bosonic full CFT one also imposes the commutativity condition
\[
  m\circ c_{A,A}=m
\]
when the algebra describes a local extension of the chiral theory.  Boundary
conditions for this full CFT are left \(A\)-modules
\[
  \rho_M:A\otimes M\to M,\qquad
  \rho_M(m\otimes\operatorname{id}_M)
  =
  \rho_M(\operatorname{id}_A\otimes\rho_M)\alpha_{A,A,M},
  \qquad
  \rho_M(\eta\otimes\operatorname{id}_M)=\operatorname{id}_M .
\]

The Frobenius identity is the algebraic form of the elementary cutting move,
not an abstract decoration.  In the ordinary finite model
\(A=\operatorname{Mat}_d(\mathbb C)\), take matrix units \(E_{ij}\), ordinary
multiplication \(m\), and trace counit \(\varepsilon(X)=\operatorname{Tr}X\).
The coproduct dual to multiplication with respect to
\(\varepsilon(XY)\) is
\begin{equation}
  \Delta(E_{ij})
  =
  \sum_{r=1}^{d} E_{ir}\otimes E_{rj}.
  \label{eq:bcft-matrix-frobenius-coproduct}
\end{equation}
Then, on \(E_{ij}\otimes E_{k\ell}\),
\begin{equation}
  (m\otimes\operatorname{id})
  (\operatorname{id}\otimes\Delta)
  =
  \Delta m
  =
  (\operatorname{id}\otimes m)
  (\Delta\otimes\operatorname{id})
  =
  \delta_{jk}\sum_{r=1}^{d} E_{ir}\otimes E_{r\ell},
  \label{eq:bcft-matrix-frobenius-cutting-move}
\end{equation}
and
\begin{equation}
  m\Delta(E_{ij})=d\,E_{ij}.
  \label{eq:bcft-matrix-frobenius-specialness}
\end{equation}
Thus sliding an \(A\)-line through a multiplication junction is exactly the
Frobenius identity, while a contractible \(A\)-loop contributes the fixed
specialness scalar.  The categorical ribbon construction replaces these
matrix units by dual bases of morphism spaces and inserts associators,
braidings, and twists; the local algebraic reason for genus-zero boundary
sewing is already visible in \eqref{eq:bcft-matrix-frobenius-cutting-move}.

This language converts the open sector into a finite-dimensional statement.
For two boundary conditions \(M,N\), the multiplicity of chiral type \(i\)
in the open Hilbert space is
\begin{equation}
  n_{MN}{}^i
  =
  \dim
  \operatorname{Hom}_{A}(M\otimes U_i,N),
  \qquad
  \Hilb_{MN}
  =
  \bigoplus_{i\in I}
  \operatorname{Hom}_{A}(M\otimes U_i,N)\otimes\mathcal H_i .
  \label{eq:bcft-module-open-sector}
\end{equation}
Here \(M\otimes U_i\) is given the left \(A\)-action
\(\rho_M\otimes\operatorname{id}_{U_i}\), with the associator understood.
Thus the annulus multiplicities are dimensions of spaces of \(A\)-linear
morphisms, and integrality is no longer a separate miracle.

Boundary OPE associativity is likewise a module-theoretic composition law.
Let
\[
  f\in\operatorname{Hom}_{A}(M\otimes U_i,N),
  \qquad
  g\in\operatorname{Hom}_{A}(N\otimes U_j,P).
\]
Then
\begin{equation}
  g\circ(f\otimes\operatorname{id}_{U_j})\circ\alpha_{M,U_i,U_j}
  \in
  \operatorname{Hom}_{A}(M\otimes(U_i\otimes U_j),P).
  \label{eq:bcft-module-ope-composition}
\end{equation}
Choose a basis of chiral fusion morphisms
\[
  e^{k,\mu}_{ij}:U_i\otimes U_j\to U_k,
  \qquad
  \bar e^{ij}_{k,\mu}:U_k\to U_i\otimes U_j,
  \qquad
  e^{k,\mu}_{ij}\bar e^{ij}_{\ell,\nu}
  =
  \delta_{k\ell}\delta_{\mu\nu}\operatorname{id}_{U_k}.
\]
Inserting the resolution
\[
  \operatorname{id}_{U_i\otimes U_j}
  =
  \sum_{k,\mu}\bar e^{ij}_{k,\mu}e^{k,\mu}_{ij}
\]
in \eqref{eq:bcft-module-ope-composition} gives the boundary OPE
coefficients in \eqref{eq:boundary-changing-ope}.  Associativity of the
boundary OPE is the statement that the two ways to compose three such
\(A\)-linear morphisms agree after changing the parenthesization of
\(U_i\otimes U_j\otimes U_k\).  The only nontrivial transformation between
the two bases is the chiral fusing isomorphism, so the module associativity
identity above becomes the Cardy--Lewellen equation
\eqref{eq:cardy-lewellen-boundary-sewing}.  The pentagon for \(\alpha\) is
the algebraic reason that the five-boundary-field sewing diagrams agree.

The diagonal Cardy solution is the special case \(A=\mathbf 1\).  A left
\(\mathbf 1\)-module is just an object of \(\mathsf C\), and the elementary
modules are \(M=U_a\).  Equation~\eqref{eq:bcft-module-open-sector} reduces
to
\[
  n_{ab}{}^i
  =
  \dim\operatorname{Hom}(U_a\otimes U_i,U_b)
  =
  N_{ai}{}^b ,
\]
which is the boundary-field multiplicity
\eqref{eq:cardy-boundary-field-multiplicity} after the same orientation
convention used in the annulus.  The fusing-matrix expression for the Cardy
boundary three-point constants is precisely the expansion of
\eqref{eq:bcft-module-ope-composition} in the basis
\(\{e_{ij}^{k,\mu}\}\).

The annulus coefficients in a non-diagonal rational boundary theory carry a
stronger constraint than integrality.  For simple \(A\)-modules \(M,N\), form
the matrices
\[
  (n_i)_M{}^N
  =
  \dim\operatorname{Hom}_{A}(M\otimes U_i,N).
\]
Taking the graded trace of a boundary interval after inserting two successive
chiral labels \(i,j\) can be computed either by first summing over an
intermediate boundary module \(N\), or by fusing \(U_i\otimes U_j\) in the
chiral category.  Module associativity therefore gives the annulus
representation identity
\begin{equation}
  \sum_N (n_i)_M{}^N (n_j)_N{}^P
  =
  \sum_k N_{ij}{}^k (n_k)_M{}^P .
  \label{eq:bcft-annulus-nimrep-identity}
\end{equation}
This is the finite algebra behind the usual statement that annulus
multiplicity matrices form a nonnegative-integer matrix representation of the
fusion ring.  It is not a separate classification theorem: it is the
decategorified shadow of the same \(A\)-module composition law that produces
\eqref{eq:bcft-module-ope-composition}.

A small pointed module example shows how this differs from the diagonal Cardy
case.  Let the chiral fusion skeleton have labels
\(G=\mathbb Z_2\times\mathbb Z_2\), with
\(U_g\otimes U_h=U_{g+h}\), and take the subgroup
\(H=\{(0,0),(1,0)\}\).  Boundary labels are the two cosets \(G/H\).  The
open multiplicity matrix attached to \(g=(g_1,g_2)\) is
\begin{equation}
  (n_g)_{xH}{}^{yH}
  =
  \begin{cases}
  1,& yH=(x+g)+H,\\
  0,& \text{otherwise}.
  \end{cases}
  \label{eq:bcft-pointed-coset-nimrep}
\end{equation}
Thus \(n_{(0,0)}=n_{(1,0)}=\mathbf 1\), while
\(n_{(0,1)}=n_{(1,1)}\) swaps the two boundary labels.  Equation
\eqref{eq:bcft-annulus-nimrep-identity} is then just the coset action law
\(n_gn_h=n_{g+h}\).  The example is finite and algebraic; realizing it as a
unitary local rational BCFT requires the corresponding modular, Frobenius,
and analytic sewing data in addition to the displayed matrices.

Bulk fields in the same algebraic construction are bimodule morphisms.  If
\(U_i^+\otimes A\otimes U_j^-\) denotes the \(A\)-bimodule obtained by using
the braiding for the left chiral label and the inverse braiding for the
right chiral label, the closed bulk multiplicity is
\begin{equation}
  Z_{ij}(A)
  =
  \dim
  \operatorname{Hom}_{A|A}
  (U_i^+\otimes A\otimes U_j^-,A).
  \label{eq:bcft-frobenius-bulk-multiplicity}
\end{equation}
The diagonal invariant is again recovered from \(A=\mathbf 1\), for which
the bimodule morphism space is one-dimensional only when \(j=i^\vee\).  The
bulk-boundary OPE is obtained by evaluating such a bimodule morphism on the
boundary module \(M\) and then projecting to
\(\operatorname{End}_{A}(M)\).  This is the categorical form of the
classifying algebra mentioned above: disk one-point functions are
simultaneous eigenvalues of the commutative algebra of bulk fields acting on
elementary boundary modules.  Formula
\eqref{eq:fusion-ring-character-identity} is the \(A=\mathbf 1\) version of
this statement.

\paragraph{Finite algebra shadow of the classifying algebra.}
The preceding categorical sentence has a simple finite algebra model which
keeps the logic visible in non-diagonal examples.  Let
\[
  A_{\rm fin}=\bigoplus_{r=1}^{R}\operatorname{Mat}_{d_r}(\mathbb C)
\]
be an ordinary finite-dimensional semisimple algebra, regarded as the
decategorified shadow of a special symmetric Frobenius algebra object and
its module category.  Its center is
\[
  Z(A_{\rm fin})
  =
  \bigoplus_{r=1}^{R}\mathbb C\,e_r ,
  \qquad
  e_r e_s=\delta_{rs}e_r,\qquad
  \sum_{r=1}^{R}e_r=1 .
\]
The simple left \(A_{\rm fin}\)-modules are the column modules
\(M_r=\mathbb C^{d_r}\).  A central element
\(z=\sum_r z_r e_r\) acts on \(M_r\) by the scalar \(z_r\), and therefore
defines a character
\begin{equation}
  \chi_r:Z(A_{\rm fin})\to\mathbb C,\qquad
  \chi_r(z)=z_r .
  \label{eq:bcft-finite-center-character}
\end{equation}
The disk identity-channel sewing equation is, in this finite model, exactly
the multiplicativity of this central character:
\begin{equation}
  \chi_r(zz')=\chi_r(z)\chi_r(z'),
  \qquad z,z'\in Z(A_{\rm fin}) .
  \label{eq:bcft-finite-center-character-sewing}
\end{equation}
This is the non-diagonal analogue of
\eqref{eq:cardy-disk-two-bulk-classifying-sewing}: the bulk field that
survives in the boundary identity channel is not a raw element of the
fusion ring but its image in the commutative classifying algebra acting on
the elementary boundary module.

Two cautions are part of the statement.  First, a noncentral element of
\(A_{\rm fin}\) does not act by a scalar on \(M_r\), so it cannot be read as
a disk one-point coefficient of the identity boundary field.  Its boundary
image is an endomorphism-valued insertion.  Second, if one replaces the
elementary boundary condition \(M_r\) by a direct sum
\(M=\oplus_r M_r^{\oplus n_r}\), the boundary identity algebra becomes
\(\operatorname{End}_{A_{\rm fin}}(M)\simeq\oplus_r\operatorname{Mat}_{n_r}(\mathbb C)\).
The scalar character \(\chi_r\) is replaced by a matrix-valued action, or by
a chosen trace after adding Chan--Paton data.  Such a trace is additive over
summands but is not a multiplicative character unless the boundary module is
simple.  Thus Cardy--Lewellen disk sewing distinguishes elementary boundary
conditions from finite direct sums even before any analytic issue about
convergence of bordered-surface correlators enters.

The disk partition function defines the Affleck--Ludwig boundary entropy.
For a Cardy boundary state in a diagonal rational CFT,
\begin{equation}
  g_a
  =
  \langle0|a\rangle
  =
  \frac{S_{a0}}{\sqrt{S_{00}}},
  \label{eq:cardy-boundary-entropy}
\end{equation}
with the vacuum Ishibashi state normalized as in
\eqref{eq:ishibashi-overlap-character}.

\begin{quotedtheorem}[Boundary entropy gradient formula]
\label{qthm:boundary-entropy-gradient-formula}
Let a unitary two-dimensional QFT on a half-cylinder have a fixed conformal
bulk and a renormalized finite-dimensional chart of boundary couplings
\(\lambda^a\).  Let \(B^a(\lambda)\) be the boundary RG vector field and
\(\phi_a(\tau)\) the renormalized boundary fields conjugate to
\(\lambda^a\).  Assume that the boundary trace of the stress tensor has the
renormalized form
\begin{equation}
  \theta(\tau)
  =
  B^a(\lambda)\phi_a(\tau)+\partial_\tau j(\tau)
  \label{eq:boundary-trace-beta-fields}
\end{equation}
and quotient the coupling chart by redundant directions generated by total
boundary derivatives.  Let \(z(\lambda,L)\) be the boundary contribution to
the cylinder partition function after subtracting the bulk extensive term and
the boundary ground-state energy, and define
\begin{equation}
  s(\lambda,L)
  =
  \left(1-L\frac{\partial}{\partial L}\right)\log z(\lambda,L).
  \label{eq:boundary-entropy-away-from-fixed-point}
\end{equation}
Under the canonical short-distance and clustering hypotheses needed to
renormalize the connected boundary two-point functions, there is a positive
symmetric bilinear form
\begin{equation}
  G_{ab}(\lambda)
  =
  \frac12
  \int_0^L\!\dd\tau_1
  \int_0^L\!\dd\tau_2\,
  \left[
    1-\cos\!\left(\frac{2\pi(\tau_1-\tau_2)}{L}\right)
  \right]
  \bigl\langle\phi_a(\tau_1)\phi_b(\tau_2)\bigr\rangle_{\!c}^{\mathrm{ren}}
  \label{eq:boundary-gradient-metric}
\end{equation}
on the quotient such that
\begin{equation}
  \partial_a s(\lambda,L)
  =
  -\,G_{ab}(\lambda)B^b(\lambda).
  \label{eq:boundary-gradient-formula}
\end{equation}
\end{quotedtheorem}

The positivity in
\eqref{eq:boundary-gradient-metric} comes from the spectral representation
of the boundary two-point function.  The finite-regulator mechanism is as
follows.  Let
\(\Phi=u^a\phi_a\) be a real boundary perturbing field after subtracting its
one-point function, and first place the boundary theory on a finite
half-cylinder with discrete positive-energy boundary Hamiltonian.  Grouping
matrix elements by a positive energy gap \(\Delta_r>0\), reflection
positivity and the KMS trace give the connected Euclidean boundary two-point
function in the form
\begin{equation}
  \langle\Phi(\tau)\Phi(0)\rangle_c
  =
  \sum_r A_r
  \frac{\ee^{-\Delta_r\tau}
        +\ee^{-\Delta_r(L-\tau)}}
       {1-\ee^{-L\Delta_r}},
  \qquad
  A_r\ge0,\qquad 0<\tau<L ,
  \label{eq:bcft-boundary-spectral-two-point}
\end{equation}
where the connected trace has removed the zero-gap diagonal part.  Put
\(\kappa=2\pi/L\).  For a single gap,
\begin{align}
  \int_0^L\!\dd\tau\,
  \bigl(1-\cos \kappa\tau\bigr)
  \frac{\ee^{-\Delta\tau}
        +\ee^{-\Delta(L-\tau)}}{1-\ee^{-L\Delta}}
  &=
  \frac{2\kappa^2}{\Delta(\Delta^2+\kappa^2)} .
  \label{eq:bcft-boundary-gradient-spectral-weight}
\end{align}
Thus the quadratic form obtained from
\eqref{eq:boundary-gradient-metric} is a sum of nonnegative matrix-element
weights
\[
  L\sum_r A_r
  \frac{2\kappa^2}{\Delta_r(\Delta_r^2+\kappa^2)} .
\]
This calculation explains the shape of the trigonometric kernel in the
gradient formula: after the connected trace removes the zero-gap diagonal
part, the kernel assigns a positive finite weight to each nonzero boundary
excitation.  In the continuum theorem, the hard analytic work is to
renormalize the coincident boundary singularities, preserve the boundary Ward
identity \eqref{eq:boundary-trace-beta-fields}, and quotient total
derivative directions so that this finite spectral positivity survives in
the renormalized coupling chart.

The monotonicity statement is the direct consequence of the gradient formula,
not an independent classification theorem.  With RG time \(t=\log L\),
\begin{equation}
  \frac{d}{dt}s(\lambda(t),L)
  =
  B^a\partial_a s
  =
  -\,B^aG_{ab}B^b
  \leq0 .
  \label{eq:boundary-g-monotonicity-from-gradient}
\end{equation}
At a boundary fixed point \(B^a=0\), the quantity \(g=\exp s\) is the
Affleck--Ludwig disk degeneracy.  Therefore a boundary flow connecting two
fixed points satisfies
\[
  g_{\mathrm{UV}}\ge g_{\mathrm{IR}} .
\]
This conclusion gives an ordering constraint on boundary fixed points reached
by a given unitary boundary RG flow.  It does not determine which conformal
boundary condition is reached, nor does it supply the annulus or
Cardy--Lewellen sewing data.

\paragraph{Mechanism of the gradient formula.}
The theorem identifies a particular renormalized susceptibility with the
differential of the finite boundary entropy.  We spell out the mechanism
because the contact-term prescription, the quotient by redundant boundary
directions, and reflection positivity are part of the mathematical content of
the statement.

Use the sign convention in which the renormalized boundary action is deformed
by
\[
  S_{\partial}(\lambda)
  =
  S_{\partial}(0)+\lambda^a\int_0^L\phi_a(\tau)\,\dd\tau .
\]
Then differentiation of the renormalized boundary partition function gives
\begin{equation}
  \partial_a\log z(\lambda,L)
  =
  -\int_0^L
  \bigl\langle\phi_a(\tau)\bigr\rangle_{\lambda,L}^{\mathrm{ren}}\,
  \dd\tau .
  \label{eq:bcft-boundary-coupling-derivative}
\end{equation}
The scale Ward identity on the half-cylinder states that, after the bulk
extensive term and the boundary ground-state energy have been subtracted, the
operation \(L\partial_L\) is represented by the renormalized insertion of the
boundary trace
\[
  \Theta=\int_0^L\theta(\tau)\,\dd\tau
\]
together with local contact counterterms supported at coincident boundary
points.  Thus \((1-L\partial_L)\partial_a\log z\) contains the integrated
two-point insertion
\(\int\!\int\langle\phi_a(\tau_1)\theta(\tau_2)\rangle_c\) together with local
counterterms.  Those contact terms must be fixed by a renormalization
prescription compatible with boundary translation invariance, reflection
positivity, and the subtraction defining \(s=(1-L\partial_L)\log z\).

With these prescriptions the Ward identity can be written in the finite form
\begin{equation}
  \partial_a s(\lambda,L)
  =
  -\frac12
  \int_0^L\!\dd\tau_1
  \int_0^L\!\dd\tau_2\,
  K_L(\tau_1-\tau_2)
  \bigl\langle
    \phi_a(\tau_1)\theta(\tau_2)
  \bigr\rangle_{\!c}^{\mathrm{ren}},
  \qquad
  K_L(\tau)=1-\cos\frac{2\pi\tau}{L}.
  \label{eq:bcft-gradient-ward-susceptibility}
\end{equation}
Equation \eqref{eq:bcft-gradient-ward-susceptibility} is the real analytic
step in the gradient formula: the kernel \(K_L\) is the circle kernel, in this
normalization, that implements the finite entropy subtraction and removes the
ambiguous zero-mode and total-derivative pieces after passing to the quotient
of boundary coupling space.  Substituting the trace identity
\eqref{eq:boundary-trace-beta-fields} gives
\[
  \bigl\langle\phi_a(\tau_1)\theta(\tau_2)\bigr\rangle_c^{\mathrm{ren}}
  =
  B^b
  \bigl\langle\phi_a(\tau_1)\phi_b(\tau_2)\bigr\rangle_c^{\mathrm{ren}}
  +
  \partial_{\tau_2}
  \bigl\langle\phi_a(\tau_1)j(\tau_2)\bigr\rangle_c^{\mathrm{ren}} .
\]
The second term is a redundant direction.  On the quotient it vanishes after
integration against the Ward kernel and after the prescribed contact
subtractions.  The first term is exactly the bilinear form
\eqref{eq:boundary-gradient-metric}, and hence
\[
  \partial_a s=-G_{ab}B^b .
\]
The positivity of \(G_{ab}\) uses boundary reflection positivity of the
renormalized theory and the removal of null redundant fields, in addition to
the trigonometric Ward kernel.  With weaker analytic control the displayed
formula records the expected Ward-identity structure rather than a proved
gradient theorem.

Finally, two-dimensional BCFT fits into the broader boundary/defect language
of Chapter~\ref{ch:global-boundaries-and-defects}.  A topological line
defect can end on a boundary only when its endpoint is a specified boundary
condition-changing operator.  Noninvertible defects, such as the Ising
Kramers--Wannier line, act on the set of Cardy boundary conditions by fusion
when the required endpoint operators exist.  Boundary anomalies are likewise
not optional decorations: a boundary partition function may be a section of
an anomaly line, and anomaly cancellation is part of the boundary datum.
The present chapter supplies the Virasoro and modular-sewing machinery
special to two-dimensional conformal boundaries; the global categorical
defect viewpoint is developed in Volume IX.

\begin{openproblem}[BCFT sewing beyond rational diagonal examples]
\label{op:bcft-nonrational-sewing}
Construct a single analytic framework that includes rational Cardy
boundaries, compact-boson continuous families, the Liouville FZZT and ZZ
states of Section~\ref{sec:bcft-liouville-boundary-states},
boundary-condition-changing operators, and defect endpoints, with convergent
sewing maps for bordered Riemann surfaces.  The framework must realize the
test-space, open spectral-measure, operator-domain, pole-prescription, and
determinant-line requirements stated above, and must prove that the rational
finite-sum construction and the Liouville direct-integral construction are
two limits of the same sewing notion rather than two unrelated calculational
recipes.
\end{openproblem}
